\exnumber{03}
\extitle{UX III - Ex 4 - Entretiens}
\exfiles{ex03/fichier Figma et PDF, nomme Name\_FirstName\_DeliverableName\_Date\_VersionNumber}
\exauthorize{Those named in the exercise}
\makeheaderfiles

\textbf{Instructions :}

\begin{itemize}
\item Create an interview guide
\item Conduct 3 interviews with teenagers
\end{itemize}

After knowing more about trends and "What" people do thanks to surveys, the main goal is to understand "why".

\textbf{Livrables :}

\begin{itemize}
\item A PDF document with an interview guide
\item A PDF document with 3 scripts 
\item An affinity diagram from your interviews (PDF or digital, on Miro or Figma for example)
\item Déposez les documents sur Github, dans un zip nommé « ProjectTitle\_Name\_FirstName », et chaque livrable ainsi : « Name\_FirstName\_DeliverableName\_Date\_VersionNumber ». Exemple : « Doe\_John\_MidFiPrototype\_05242024\_V1 ».
\item Ajoutez un titre dans tous les documents écrits.
\item Soignez l'aspect global des livrables. Ils doivent être clairs, simples et lisibles. Vous pouvez regarder canva.com, sans surcharger le dessin de détails.
\end{itemize}

\textbf{Ressources :}

\begin{itemize}
\item Respondent - Tool to find users \url{https://www.respondent.io}
\item How to conduct user interviews \url{https://uxdesign.cc/how-to-conduct-user-interviews-fe4b8c34b0b7?sk=9625632a522a9309a91297191e668565}
\item How to design better products through user interviews \url{https://uxdesign.cc/how-to-design-better-products-through-user-interviews-4c5142bb1fc4}
\item Asking the right questions \url{https://uxdesign.cc/asking-the-right-questions-on-user-research-interviews-and-testing-427261742a67}
\end{itemize}

\textbf{Pistes :}

Voici les étapes recommandées :

1. Define your objective, what you are looking for when interviewing people. Usually, this has to do with your user's habits, and the most painful problems they have regarding this topic.
2. Write down 1 to 3 hypotheses you have regarding this topic, that needs to be either confirmed or disproved.
3. Choisissez 3 à 5 sujets sur lesquels il vous faut de l'information.
4. For each topic, write 1 to 3 open-ended questions. Remember, your user cannot answer with "Yes" or "No"! Ex: How, Why, What, When, Tell me about the last time you did... etc. If you really need to introduce a new topic by a close-ended question, ask why right after.
5. Write some screening questions to have a few pieces of demographics (age, job, location, etc.)
6. Structure your questions in a funnel: start with broad questions to focus on specific topics or situation. Remember to stay neutral in your tone to let the user tell their own story. Try to focus on 10 main questions.
7. Write a greeting introduction, describing briefly (1) who you are and (2) what you are looking for. No need to say too much, but enough to provide a framework and the purpose of the interview.
8. Write a conclusion to thank your user for their time, and prepare the next steps. Ex: ask to be put in touch with someone they know and who could be interested in this topic, etc.

